% Options for packages loaded elsewhere
\PassOptionsToPackage{unicode}{hyperref}
\PassOptionsToPackage{hyphens}{url}
\documentclass[
  12pt,
]{ctexart}
\usepackage{xcolor}
\usepackage{amsmath,amssymb}
\setcounter{secnumdepth}{-\maxdimen} % remove section numbering
\usepackage{iftex}
\ifPDFTeX
  \usepackage[T1]{fontenc}
  \usepackage[utf8]{inputenc}
  \usepackage{textcomp} % provide euro and other symbols
\else % if luatex or xetex
  \usepackage{unicode-math} % this also loads fontspec
  \defaultfontfeatures{Scale=MatchLowercase}
  \defaultfontfeatures[\rmfamily]{Ligatures=TeX,Scale=1}
\fi
\usepackage{lmodern}
\ifPDFTeX\else
  % xetex/luatex font selection
\fi
% Use upquote if available, for straight quotes in verbatim environments
\IfFileExists{upquote.sty}{\usepackage{upquote}}{}
\IfFileExists{microtype.sty}{% use microtype if available
  \usepackage[]{microtype}
  \UseMicrotypeSet[protrusion]{basicmath} % disable protrusion for tt fonts
}{}
\makeatletter
\@ifundefined{KOMAClassName}{% if non-KOMA class
  \IfFileExists{parskip.sty}{%
    \usepackage{parskip}
  }{% else
    \setlength{\parindent}{0pt}
    \setlength{\parskip}{6pt plus 2pt minus 1pt}}
}{% if KOMA class
  \KOMAoptions{parskip=half}}
\makeatother
\usepackage{longtable,booktabs,array}
\usepackage{calc} % for calculating minipage widths
% Correct order of tables after \paragraph or \subparagraph
\usepackage{etoolbox}
\makeatletter
\patchcmd\longtable{\par}{\if@noskipsec\mbox{}\fi\par}{}{}
\makeatother
% Allow footnotes in longtable head/foot
\IfFileExists{footnotehyper.sty}{\usepackage{footnotehyper}}{\usepackage{footnote}}
\makesavenoteenv{longtable}
\setlength{\emergencystretch}{3em} % prevent overfull lines
\providecommand{\tightlist}{%
  \setlength{\itemsep}{0pt}\setlength{\parskip}{0pt}}
\usepackage{bookmark}
\IfFileExists{xurl.sty}{\usepackage{xurl}}{} % add URL line breaks if available
\urlstyle{same}
\hypersetup{
  hidelinks,
  pdfcreator={LaTeX via pandoc}}

\author{SCX}
\date{2026}

\begin{document}

\section{SCX Mathematical Theory: Unified Theorem
Document}\label{scx-mathematical-theory-unified-theorem-document}

\begin{quote}
\textbf{Version}: 2026-06-28 \textbar{} \textbf{Status}: Complete
\textbar{} \textbf{Audit}: Post-verification \textbf{Purpose}: Single
source of truth for all 6 theorems in the SCX theory chain, for Nature
SI submission. \textbf{Verification status}: All formulas checked for
sign errors, all inequality directions verified, all constants traced to
definitions, all cross-file notation conflicts resolved.
\end{quote}

\begin{center}\rule{0.5\linewidth}{0.5pt}\end{center}

\subsection{Table of Contents}\label{table-of-contents}

\begin{enumerate}
\def\labelenumi{\arabic{enumi}.}
\setcounter{enumi}{-1}
\tightlist
\item
  \hyperref[0-preface]{Preface}

  \begin{itemize}
  \tightlist
  \item
    0.1 \hyperref[01-unified-notation-glossary]{Unified Notation
    Glossary}
  \item
    0.2 \hyperref[02-assumption-catalog-a1-a6]{Assumption Catalog
    (A1-A6)}
  \item
    0.3 \hyperref[03-dependency-graph]{Dependency Graph}
  \item
    0.4 \hyperref[04-k-disambiguation]{K Disambiguation}
  \end{itemize}
\item
  \hyperref[1-theorem-1-noise-detection-guarantee]{Theorem 1: Noise
  Detection Guarantee}
\item
  \hyperref[2-theorem-2-weak-feature-failure-lower-bound]{Theorem 2:
  Weak Feature Failure Lower Bound}
\item
  \hyperref[3-theorem-3-noise-difficulty-unidentifiability]{Theorem 3:
  Noise-Difficulty Unidentifiability}
\item
  \hyperref[4-theorem-4-exact-constant-minimax-optimality]{Theorem 4':
  Exact Constant Minimax Optimality}
\item
  \hyperref[5-theorem-5-fixed-k-cluster-consistency]{Theorem 5: Fixed-K
  Cluster Consistency}
\item
  \hyperref[6-proposition-6-bootstrap-stability-diagnostic-for-feature-strength]{Proposition
  6: Bootstrap Stability Diagnostic for Feature Strength}
\item
  \hyperref[7-cross-theorem-consistency-audit]{Cross-Theorem Consistency
  Audit}
\item
  \hyperref[8-known-gaps-and-limitations]{Known Gaps and Limitations}
\item
  \hyperref[9-si-writing-guide]{SI Writing Guide}
\item
  \hyperref[10-references]{References}
\end{enumerate}

\begin{center}\rule{0.5\linewidth}{0.5pt}\end{center}

\subsection{0. Preface}\label{preface}

\subsubsection{0.1 Unified Notation
Glossary}\label{unified-notation-glossary}

All symbols used across the 6 theorems are collected here. Conflicting
uses are disambiguated with subscripts.

\begin{longtable}[]{@{}
  >{\raggedright\arraybackslash}p{(\linewidth - 6\tabcolsep) * \real{0.2424}}
  >{\raggedright\arraybackslash}p{(\linewidth - 6\tabcolsep) * \real{0.2727}}
  >{\raggedright\arraybackslash}p{(\linewidth - 6\tabcolsep) * \real{0.2727}}
  >{\raggedright\arraybackslash}p{(\linewidth - 6\tabcolsep) * \real{0.2121}}@{}}
\toprule\noalign{}
\begin{minipage}[b]{\linewidth}\raggedright
Symbol
\end{minipage} & \begin{minipage}[b]{\linewidth}\raggedright
Meaning
\end{minipage} & \begin{minipage}[b]{\linewidth}\raggedright
Used in
\end{minipage} & \begin{minipage}[b]{\linewidth}\raggedright
Notes
\end{minipage} \\
\midrule\noalign{}
\endhead
\bottomrule\noalign{}
\endlastfoot
\(\mathcal{X}\) & Input space & Thm 1,2,3,5 & \\
\(\mathcal{Y}\) & Label space & Thm 1,3 &
\(|\mathcal{Y}| = K_{\mathcal{Y}}\) \\
\(K_{\mathcal{Y}}\) & Number of classes & Thm 1,3 & Disambiguated from
\(K_{\mathcal{S}}\) \\
\(\mathcal{S}\) & State space (index set) & Thm 1,2,3 &
\(|\mathcal{S}| = K_{\mathcal{S}}\) \\
\(K_{\mathcal{S}}\) & Number of states & Thm 2,5 & Disambiguated from
\(K_{\mathcal{Y}}\) \\
\(K\) & Generic state/class count & \emph{various} & \textbf{Always
disambiguated in this document} \\
\(s(x) : \mathcal{X} \to \mathcal{S}\) & True state assignment & Thm
1,2,3,5 & Unobserved during clustering \\
\(\hat{s}(x)\) & Estimated state assignment & Thm 2,5 & From k-means on
\(\phi(x)\) \\
\(\Pi = \{s_1,\dots,s_{K_{\mathcal{S}}}\}\) & State partition of
\(\mathcal{X}\) & Thm 1 & Measurable partition \\
\(\rho_s = \mathbb{P}(X \in s)\) & State probability & Thm 1 &
\(\sum_s \rho_s = 1\) \\
\(\mu_s\) & State-\(s\) clean error upper bound & Thm 1 &
\(\sup_{x\in s} \mathbb{E}[C \mid \text{clean}, X=x] \leq \mu_s\) \\
\(C_{\text{bal}}\) & Error balance constant (\(\geq 1\)) & Thm 1 &
Controls error concentration \\
\(f^* : \mathcal{X} \to \mathcal{Y}\) & True oracle & Thm 1,3 &
Unobserved \\
\(\{f_m\}_{m=1}^M\) & Expert models & Thm 1,2,3 & \(M\) experts \\
\(e_m(x,y) = \mathbf{1}\{\ell(f_m(x),y) > \tau\}\) & Expert error
indicator & Thm 1 & \\
\(C(x) = \frac{1}{M}\sum_m e_m(x,y)\) & Consensus score & Thm 1 & \\
\(\theta\) & Detection threshold & Thm 1,4' & \\
\(\eta\) & Global noise rate & Thm 1,2,3,4' & \(\eta \in (0, 1/2)\)
typically \\
\(\ell : \mathcal{Y} \times \mathcal{Y} \to [0,B]\) & Bounded loss & Thm
1 & \\
\(\tau > 0\) & Expert error threshold & Thm 1 & \\
\(\Delta_s\) & State-\(s\) separation gap & Thm 1 & \\
\(\phi(X)\) & Feature representation & Thm 2,5 &
\(\phi: \mathcal{X} \to \mathbb{R}^{d_\phi}\) \\
\(\delta = I(\phi(X); S)\) & Feature-state mutual info & Thm 2 & In
nats \\
\(\varepsilon_\phi = \delta / \log K_{\mathcal{S}}\) & Normalized
feature weakness & Thm 2 & \\
\(Z \in \{0,1\}\) & Noise indicator & Thm 2 & \(Z = 1\) for noise \\
\(D_m(x) = \ell(f_m(x), y)\) & Expert loss & Thm 2 & \\
\(NS(x)\) & SCX noise score & Thm 2 & \\
\(C_F\) & F1 Lipschitz constant & Thm 2 & \(C_F \leq 2\) in operating
range \\
\(p_0 = \mu_s\) & Clean error rate (state \(s\)) & Thm 4' & \\
\(p_1 = 1 - C_{\text{bal}} \cdot \mu_s/(K_{\mathcal{Y}}-1)\) & Noise
error rate (state \(s\)) & Thm 4' & \\
\(\kappa = C(\text{Bern}(p_0), \text{Bern}(p_1))\) & Chernoff
information & Thm 4' &
\(\kappa = \text{KL}(\theta^*\|p_0) = \text{KL}(\theta^*\|p_1)\) \\
\(\theta^*\) & Chernoff point & Thm 4' & \(\theta^* \in (p_0, p_1)\) \\
\(\lambda_0^*, \lambda_1^*\) & Saddlepoints at \(\theta^*\) & Thm 4' &
\(\lambda_0^* > 0, \lambda_1^* < 0\) \\
\(D^* = \lambda_0^* + |\lambda_1^*|\) & Total log-odds & Thm 4' & \\
\(s = |\lambda_1^*| / D^*\) & Exponent fraction & Thm 4' &
\(s \in (0,1)\) \\
\(\theta^\dagger\) (or \(\theta_{\text{opt}}\)) & Adaptive threshold &
Thm 4' & \(\theta^\dagger = \theta^* + O(1/M)\) \\
\(C_{\min}\) & Minimax optimal constant & Thm 4' & Lemma E canonical
form \\
\(\mu_k\) & True cluster center (state \(k\)) & Thm 5 & \\
\(\Delta_{\min}\) & Minimum state separation & Thm 5 &
\(\min_{i\neq j} \|\mu_i - \mu_j\|_2\) \\
\(\sigma^2\) & Sub-Gaussian variance proxy & Thm 5 & \\
\(d_\phi\) & Feature dimension & Thm 5 & \\
\(n_{\min}\) & Minimum per-state sample size & Thm 5 & \\
\(S(\Phi, K)\) & Clustering stability score & Prop 6 & Mean ARI over
bootstrap \\
\(\mathcal{P}_{\text{noise}}\) & Noise-world distribution & Thm 3 & \\
\(\mathcal{P}_{\text{hard}}\) & Hardness-world distribution & Thm 3 & \\
\end{longtable}

\subsubsection{0.2 Assumption Catalog
(A1-A6)}\label{assumption-catalog-a1-a6}

\textbf{All 6 theorems share the SCX assumption framework}, though
individual theorems only require subsets.

\begin{longtable}[]{@{}
  >{\raggedright\arraybackslash}p{(\linewidth - 8\tabcolsep) * \real{0.0638}}
  >{\raggedright\arraybackslash}p{(\linewidth - 8\tabcolsep) * \real{0.2340}}
  >{\raggedright\arraybackslash}p{(\linewidth - 8\tabcolsep) * \real{0.3617}}
  >{\raggedright\arraybackslash}p{(\linewidth - 8\tabcolsep) * \real{0.1915}}
  >{\raggedright\arraybackslash}p{(\linewidth - 8\tabcolsep) * \real{0.1489}}@{}}
\toprule\noalign{}
\begin{minipage}[b]{\linewidth}\raggedright
\#
\end{minipage} & \begin{minipage}[b]{\linewidth}\raggedright
Assumption
\end{minipage} & \begin{minipage}[b]{\linewidth}\raggedright
Formal Statement
\end{minipage} & \begin{minipage}[b]{\linewidth}\raggedright
Used by
\end{minipage} & \begin{minipage}[b]{\linewidth}\raggedright
Notes
\end{minipage} \\
\midrule\noalign{}
\endhead
\bottomrule\noalign{}
\endlastfoot
\textbf{A1} & Disjoint Training Sets & \(D_m \sim \mathcal{D}^{n_m}\),
\(D_m \cap D_{m'} = \varnothing\), \(D_m \perp D_{m'}\) & Thm 1, Thm 3 &
Essential for expert independence \\
\textbf{A2} & Conditional Independence (Clean) &
\(\{e_m(x,y)\}_{m=1}^M\) conditionally independent given \(x\) for clean
samples & Thm 1, Thm 3 & Follows from A1 \\
\textbf{A3} & Bounded Loss & \(\ell(a,b) \in [0, B]\), \(B < \infty\) &
Thm 1 & Technical (Hoeffding/Chernoff) \\
\textbf{A4} & Uniform Independent Noise & Noise events \(\perp x\),
\(\perp D_m\); noise label uniform over \(\mathcal{Y}\setminus\{y^*\}\)
& Thm 1, Thm 3 & Breaks unidentifiability \\
\textbf{A5} & State Homogeneity &
\(\sup_{x\in s} \mathbb{E}[C \mid \text{clean}, X=x] \leq \mu_s\) per
state \(s\) & Thm 1, Thm 4' & Within-state error uniformity \\
\textbf{A6} & Balanced Error Distribution &
\(\max_{c\neq y^*} \mu_c(x) \leq C_{\text{bal}} \cdot \mu_s/(K_{\mathcal{Y}}-1)\)
& Thm 1 & Breaks unidentifiability when \(C_{\text{bal}}=1\) \\
\end{longtable}

\textbf{Minimal subsets that break Theorem 3's unidentifiability} (see
Theorem 3, Corollaries 2-5): - \(\{A1, A4, A5\}\) (training independence
+ uniform noise + state homogeneity) - \(\{A1, A4, A6\}\) (training
independence + uniform noise + balanced errors) - \(\{A5, A6\}\) with
\(|\mathcal{S}| \geq 2\) (state homogeneity + balanced errors + multiple
states)

\subsubsection{0.3 Dependency Graph}\label{dependency-graph}

\begin{verbatim}
Theorem 3 (Unidentifiability)   ──── provides necessity justification for ────→ A1-A6
                                                                                      │
                                                                                      ▼
Theorem 1 (Noise Detection)   ──── positive result under A1-A6 ────→  F1 ≥ 1 - O(e^{-2MΔ²})
        │                                                                          
        │ provides error model (Bernoulli(p₀), Bern(p₁))                           
        ▼                                                                          
Theorem 4' (Exact Constant)   ──── refines Theorem 1's rate ────→  exact constant κ = KL(θ*‖p₀)
        │                                                                          
        │ provides expert error model                                               
        ▼                                                                          
Theorem 2 (Weak Feature)      ──── limits when features δ-weak ────→  F1_SCX ≤ F1_base + O(√δ)
        │                                                                          
        │ positive counterpart: strong features ⇒ state discovery succeeds         
        ▼                                                                          
Theorem 5 (Cluster Consistency) ──── recovery of true partition ──→  P(error) ≤ K·exp(-c·n_min·Δ²_min)
        │                                                                          
        │ practical diagnostic for Theorem 2's δ                                    
        ▼                                                                          
Proposition 6 (Stability Diag.) ──── bootstrap stability test ────→  S(Φ,K) > 0.7 ⇒ features not bottleneck
\end{verbatim}

\subsubsection{0.4 K Disambiguation}\label{k-disambiguation}

\textbf{CRITICAL: The symbol \(K\) is overloaded across files.} This
document uses:

\begin{longtable}[]{@{}
  >{\raggedright\arraybackslash}p{(\linewidth - 6\tabcolsep) * \real{0.2432}}
  >{\raggedright\arraybackslash}p{(\linewidth - 6\tabcolsep) * \real{0.2703}}
  >{\raggedright\arraybackslash}p{(\linewidth - 6\tabcolsep) * \real{0.2432}}
  >{\raggedright\arraybackslash}p{(\linewidth - 6\tabcolsep) * \real{0.2432}}@{}}
\toprule\noalign{}
\begin{minipage}[b]{\linewidth}\raggedright
Context
\end{minipage} & \begin{minipage}[b]{\linewidth}\raggedright
Notation
\end{minipage} & \begin{minipage}[b]{\linewidth}\raggedright
Meaning
\end{minipage} & \begin{minipage}[b]{\linewidth}\raggedright
Used in
\end{minipage} \\
\midrule\noalign{}
\endhead
\bottomrule\noalign{}
\endlastfoot
Classification & \(K_{\mathcal{Y}}\) & Number of classes/labels,
\(K_{\mathcal{Y}} = |\mathcal{Y}|\) & Thm 1, 3, 4' \\
State discovery & \(K_{\mathcal{S}}\) & Number of states,
\(K_{\mathcal{S}} = |\mathcal{S}|\) & Thm 2, 5, Prop 6 \\
\end{longtable}

When reusing notation from source files, the source file's \(K\) is
mapped as follows: - \textbf{Theorem 1}: \(K\) in source =
\(K_{\mathcal{Y}}\) (number of classes). Source \(\mathcal{S}\) is the
state space. - \textbf{Theorem 2}: \(K\) in source = \(K_{\mathcal{S}}\)
(number of states). - \textbf{Theorem 3}: \(K\) in source =
\(K_{\mathcal{Y}}\) (number of classes). - \textbf{Theorem 4'}: \(K\) in
source = \(K_{\mathcal{Y}}\) (number of classes, appears only in \(p_1\)
formula). - \textbf{Theorem 5}: \(K\) in source = \(K_{\mathcal{S}}\)
(number of states/clusters). - \textbf{Proposition 6}: \(K\) in source =
\(K_{\mathcal{S}}\) (number of clusters).

\begin{center}\rule{0.5\linewidth}{0.5pt}\end{center}

\subsection{1. Theorem 1: Noise Detection
Guarantee}\label{theorem-1-noise-detection-guarantee}

\subsubsection{1.1 Polished Statement}\label{polished-statement}

\textbf{Theorem 1 (SCX Noise Detection Guarantee).} Let assumptions
(A1)-(A6) hold. Let \(\rho_s = \mathbb{P}(X \in s)\) be the state
probability for \(s \in \mathcal{S}\). For any threshold \(\theta\)
satisfying, for a given state \(s\),

\[\mu_s < \theta < 1 - C_{\text{bal}} \cdot \frac{\mu_s}{K_{\mathcal{Y}}-1},\]

define the state-level separation gap

\[\Delta_s = \min\!\left(\theta - \mu_s,\; 1 - C_{\text{bal}} \cdot \frac{\mu_s}{K_{\mathcal{Y}}-1} - \theta\right) > 0.\]

Then the SCX noise detector achieves the following F1 lower bound:

\[\boxed{\;\text{F1} \;\geq\; 1 - \frac{1}{\eta} \sum_{s \in \mathcal{S}} \rho_s \cdot \exp\!\bigl(-2M\Delta_s^2\bigr)\;}\]

Equivalently, using the tighter Chernoff form:

\[\text{F1} \;\geq\; 1 - \frac{1}{\eta} \sum_{s \in \mathcal{S}} \rho_s \cdot \Bigl[ \exp\!\bigl(-M \cdot \text{KL}(\theta \,\|\, \mu_s)\bigr) + \frac{1-\eta}{\eta} \cdot \exp\!\bigl(-M \cdot \text{KL}(\theta \,\|\, 1 - C_{\text{bal}} \cdot \frac{\mu_s}{K_{\mathcal{Y}}-1})\bigr) \Bigr]\]

\textbf{Asymptotic corollary}: As \(M \to \infty\), for all states with
\(\mu_s < (K_{\mathcal{Y}}-1)/K_{\mathcal{Y}}\),

\[\text{F1} = 1 - \mathcal{O}_P\!\left(\frac{1}{\eta} \cdot e^{-2M\Delta_{\min}^2}\right), \quad \Delta_{\min} = \min_{s \in \mathcal{S}} \Delta_s.\]

\subsubsection{1.2 Key Lemmas (Compressed)}\label{key-lemmas-compressed}

\textbf{Lemma 1 (Mean Separation).} Under (A1)-(A5), for any
\(x \in s\):

\[\mathbb{E}[C \mid \text{clean}, X = x] \leq \mu_s\]
\[\mathbb{E}[C \mid \text{noise}, X = x] = 1 - \frac{1}{K_{\mathcal{Y}}-1} \cdot \mathbb{E}[C \mid \text{clean}, X = x] \geq 1 - \frac{\mu_s}{K_{\mathcal{Y}}-1}\]

The gap
\(\mathbb{E}[C \mid \text{noise}] - \mathbb{E}[C \mid \text{clean}] \geq 1 - \frac{K_{\mathcal{Y}}}{K_{\mathcal{Y}}-1}\mu_s\)
is positive when \(\mu_s < (K_{\mathcal{Y}}-1)/K_{\mathcal{Y}}\).

\textbf{Lemma 2 (FPR Upper Bound).} For states with \(\mu_s < \theta\):

\[\mathbb{P}(C > \theta \mid \text{clean}, X \in s) \leq \exp\!\bigl(-2M(\theta - \mu_s)^2\bigr)\]

\emph{Proof}: Hoeffding inequality on conditionally independent
\(\{e_m\}\), using \(\mathbb{E}[C \mid \text{clean}] \leq \mu_s\) and
\(e_m \in [0,1]\).

\textbf{Lemma 3 (TPR Lower Bound).} Under (A1)-(A6), for states with
\(\theta < 1 - C_{\text{bal}} \cdot \mu_s/(K_{\mathcal{Y}}-1)\):

\[\mathbb{P}(C > \theta \mid \text{noise}, X \in s) \geq 1 - \exp\!\left(-2M\!\left(1 - C_{\text{bal}} \cdot \frac{\mu_s}{K_{\mathcal{Y}}-1} - \theta\right)^2\right)\]

\emph{Proof}: Condition on noise label \(c\), apply Hoeffding to each
conditional Bernoulli (using A6 to bound
\(\mathbb{E}[C \mid x, c] \geq 1 - C_{\text{bal}}\cdot\mu_s/(K_{\mathcal{Y}}-1)\)),
then average over \(c\).

\subsubsection{1.3 Verification Checks}\label{verification-checks}

\textbf{Check: Hoeffding applications correct.} - Lemma 2:
\(C = \frac{1}{M}\sum e_m\) with \(e_m \in [0,1]\),
\(\mathbb{E}[C] \leq \mu_s\). For \(\theta > \mu_s\),
\(\mathbb{P}(C - \mathbb{E}[C] > \theta - \mu_s) \leq \exp(-2M(\theta-\mu_s)^2)\).
Direction: upper tail (error above threshold). Correct. \(\checkmark\) -
Lemma 3:
\(\mathbb{E}[C \mid x,c] \geq 1 - C_{\text{bal}}\cdot\mu_s/(K_{\mathcal{Y}}-1)\).
For \(\theta < 1 - C_{\text{bal}}\cdot\mu_s/(K_{\mathcal{Y}}-1)\),
\(\mathbb{P}(\mathbb{E}[C] - C > \mathbb{E}[C] - \theta) \leq \exp(-2M(\mathbb{E}[C]-\theta)^2)\).
Direction: lower tail (consensus below threshold despite it being
noise). The gap condition uses \(\mathbb{E}[C] - \theta > 0\). Correct.
\(\checkmark\)

\textbf{Check: \(C_{\text{bal}}\) used consistently.} - Lemma 1 (mean):
Does NOT use \(C_{\text{bal}}\) because the mean expression
\(\mathbb{E}[C \mid \text{noise}, x] = 1 - \frac{1}{K_{\mathcal{Y}}-1} \mathbb{E}[C \mid \text{clean}, x]\)
holds regardless of error concentration --- the uniform noise label (A4)
averages over all \(c \neq y^*\) regardless of how errors distribute
across classes. Correct. \(\checkmark\) - Lemma 3 (concentration): DOES
use \(C_{\text{bal}}\) because when conditioning on a specific noise
label \(c\), the expert error probability is
\(\mathbb{P}(f_m(x) \neq c \mid x) \geq 1 - C_{\text{bal}}\cdot\mu_s/(K_{\mathcal{Y}}-1)\)
by A6. Without A6, concentration could fail for the worst-case \(c\).
Correct. \(\checkmark\) - Theorem 1 statement: \(\Delta_s\) uses
\(C_{\text{bal}}\) in the noise-side term. The original definition
(\(C_{\text{bal}}=1\)) is recovered as a special case. Correct.
\(\checkmark\)

\textbf{Check: F1 bound formula matches proof.} - Start:
\(\text{TPR} \geq 1 - \delta_1\), \(\text{FPR} \leq \delta_2\) where
\(\delta_1 = \sum_s \rho_s \cdot \exp(-2M(1 - C_{\text{bal}}\cdot\mu_s/(K_{\mathcal{Y}}-1) - \theta)^2)\),
\(\delta_2 = \sum_s \rho_s \cdot \exp(-2M(\theta - \mu_s)^2)\). - F1
expression:
\(\frac{2\eta \cdot \text{TPR}}{\eta(1+\text{TPR}) + (1-\eta)\text{FPR}} \geq \frac{2\eta(1-\delta_1)}{\eta(2-\delta_1) + (1-\eta)\delta_2}\).
- Bound manipulation:
\(1 - \text{F1} \leq \delta_1 + \frac{1-\eta}{\eta}\delta_2\) (verified
algebraically in Lemma B). - Using
\(\Delta_s = \min(\theta - \mu_s, 1 - C_{\text{bal}}\cdot\mu_s/(K_{\mathcal{Y}}-1) - \theta)\),
we have \(\exp(-2M(\theta - \mu_s)^2) \leq \exp(-2M\Delta_s^2)\) and
similarly for the noise term. Therefore:
\[\text{F1} \geq 1 - \sum_s \rho_s[\exp(-2M\Delta_s^2) + \frac{1-\eta}{\eta}\exp(-2M\Delta_s^2)] = 1 - \frac{1}{\eta}\sum_s \rho_s \exp(-2M\Delta_s^2).\]
Correct. \(\checkmark\)

\textbf{Check: Chernoff bound KL direction.} - For the noise lower tail:
\(\mathbb{P}(C \leq \theta \mid \text{noise}) \leq \exp(-M \cdot \text{KL}(\theta \,\|\, 1 - C_{\text{bal}}\cdot\mu_s/(K_{\mathcal{Y}}-1)))\).
- Standard Chernoff: for i.i.d. Bernoulli with mean \(p\),
\(\mathbb{P}(\bar{X} \leq \theta) \leq \exp(-n \cdot \text{KL}(\theta \| p))\)
when \(\theta < p\). Here
\(p = 1 - C_{\text{bal}}\cdot\mu_s/(K_{\mathcal{Y}}-1)\), \(\theta < p\)
by assumption. So KL direction is \(\text{KL}(\theta \| p)\) with first
arg = threshold, second = true mean. Correct. \(\checkmark\) -
\textbf{Historical note}: The 2026-06-27 correction changed
\(\text{KL}(1-\theta \| 1 - \mu_s/(K_{\mathcal{Y}}-1))\) to
\(\text{KL}(\theta \| 1 - C_{\text{bal}}\cdot\mu_s/(K_{\mathcal{Y}}-1))\),
which is correct. \(\checkmark\)

\begin{center}\rule{0.5\linewidth}{0.5pt}\end{center}

\subsection{2. Theorem 2: Weak Feature Failure Lower
Bound}\label{theorem-2-weak-feature-failure-lower-bound}

\subsubsection{2.1 Polished Statement}\label{polished-statement-1}

\textbf{Theorem 2 (Weak Feature Failure Lower Bound).} Let
\(\phi: \mathcal{X} \to \mathbb{R}^{d_\phi}\) be a \(\delta\)-weak
feature mapping with respect to the true state \(S\), i.e.,
\(I(\phi(X); S) \leq \delta\) (in nats). Let \(h_{\text{SCX}}\) be the
SCX noise detector operating under the standard pipeline (clustering on
\(\phi\), state-conditional reliability estimates, noise scoring).
Assume the estimated states are approximately balanced
(\(\max_{\hat{s}} \rho(\hat{s}) / \min_{\hat{s}} \rho(\hat{s}) \leq R\))
and the clustering algorithm does not introduce error beyond the Fano
lower bound. Then:

\textbf{(a) AUC bound:}
\[\boxed{\;AUC(h_{\text{SCX}}) \leq AUC_{\text{base}} + \sqrt{\frac{\delta}{2}} \cdot \left(\frac{1}{\eta} + \frac{1}{1-\eta}\right)\;}\]

\textbf{(b) PR-AUC bound:}
\[\boxed{\;PRAUC(h_{\text{SCX}}) \leq PRAUC_{\text{base}} + \sqrt{\frac{\delta}{2}} \cdot \left(\frac{1}{\eta} + \frac{1}{1-\eta}\right)\;}\]

\textbf{(c) F1 bound:} There exists a constant \(C_F\) (depending on the
minimum precision/recall; \(C_F \leq 2\) in the typical operating range
where precision, recall \(\geq 0.1\)) such that:
\[\boxed{\;F1(h_{\text{SCX}}) \leq F1_{\text{base}} + C_F \cdot \sqrt{\frac{\delta}{2}}\;}\]

Here \(AUC_{\text{base}}\), \(PRAUC_{\text{base}}\),
\(F1_{\text{base}}\) are the performance metrics of the loss-based
baseline detector that thresholds \(\max_m \ell(f_m(x), y)\) without
using \(\phi\) or state information.

\subsubsection{2.2 Verification Checks}\label{verification-checks-1}

\textbf{Check: \(\sqrt{\delta/2}\) used consistently everywhere.} -
Pinsker:
\(TV(P, \tilde{P}) \leq \sqrt{KL(P\|\tilde{P})/2} = \sqrt{\delta/2}\).
The \(\tilde{P}\) distribution forces \(\phi \perp S\) while preserving
marginals. Correct. \(\checkmark\) - Conditional TV:
\(TV(P(\cdot|Z=1), \tilde{P}(\cdot|Z=1)) \leq TV(P, \tilde{P}) / \eta \leq \frac{1}{\eta}\sqrt{\delta/2}\).
Derived via
\(|P(A|Z=1) - \tilde{P}(A|Z=1)| = |P(A\cap\{Z=1\}) - \tilde{P}(A\cap\{Z=1\})|/\eta\).
Since \(\tilde{P}\) preserves \(P(Z=1) = \eta\), denominator is correct.
\(\checkmark\) - AUC:
\(|AUC_P - AUC_{\tilde{P}}| \leq TV(P(\cdot|Z=1), \tilde{P}(\cdot|Z=1)) + TV(P(\cdot|Z=0), \tilde{P}(\cdot|Z=0))\)
because AUC = \(\mathbb{E}[\mathbf{1}\{s_n > s_c\}]\) with
\(s_n \perp s_c\) (independent sampling), and
\(TV(\text{product}) \leq TV(\text{marginal}_1) + TV(\text{marginal}_2)\).
The \(\sqrt{\delta/2}\) factor propagates correctly. \(\checkmark\) -
F1:
\(|F1_P - F1_{\tilde{P}}| \leq C_F \cdot TV(P_{\text{pred}}, \tilde{P}_{\text{pred}}) \leq C_F \cdot \sqrt{\delta/2}\).
F1 is a function of the joint distribution \(P(\hat{z}, Z)\), not
conditionals, so no \(\eta\) amplification. \(\checkmark\)

\textbf{Check: Pinsker \(\to\) TV \(\to\) F1 chain correct.} - Step 1:
Construct \(\tilde{P}\) with \(\tilde{P}(\phi, S) = P(\phi)P(S)\).
\(KL(P\|\tilde{P}) = I(\phi; S) = \delta\). \(\checkmark\) - Step 2: By
data processing inequality,
\(TV(P_{\text{pred}}, \tilde{P}_{\text{pred}}) \leq TV(P, \tilde{P})\)
where \(P_{\text{pred}}\) is the joint of
\((\hat{z}_{\text{SCX}}(X), Z)\). \(\checkmark\) - Step 3: Under
\(\tilde{P}\), \(\phi \perp S\), so SCX state discovery fails (Lemma 2
of source: state consistency estimates converge to global mean
\(\bar{C}\), noise score proportional to loss). Thus SCX detector
\(\equiv\) loss baseline under \(\tilde{P}\). \(\checkmark\) - Step 4:
By Pinsker, \(TV(P, \tilde{P}) \leq \sqrt{\delta/2}\). So
\(|F1_P - F1_{\tilde{P}}| \leq C_F \cdot \sqrt{\delta/2}\). Since
\(F1_{\tilde{P}} = F1_{\text{base}}\), we get
\(F1_P \leq F1_{\text{base}} + C_F\sqrt{\delta/2}\). \(\checkmark\) -
\textbf{Sign check}: The inequality is
\(F1_P \leq F1_{\tilde{P}} + C_F \cdot TV(P, \tilde{P})\), which is an
upper bound. This is correct --- the theorem says SCX cannot exceed
baseline by more than \(C_F\sqrt{\delta/2}\), not that it's
lower-bounded. \(\checkmark\)

\textbf{Check: AUC/PR-AUC \(\eta\)-dependent bounds.} -
\(|AUC_P - AUC_{\tilde{P}}| \leq \frac{1}{\eta}\sqrt{\frac{\delta}{2}} + \frac{1}{1-\eta}\sqrt{\frac{\delta}{2}} = \sqrt{\frac{\delta}{2}}(\frac{1}{\eta} + \frac{1}{1-\eta})\).
- This bound becomes loose when \(\eta \to 0\) or \(\eta \to 1\). The
document correctly acknowledges this:
``\(\eta\)越小，界越宽松------这反映了检测稀有噪声的内在困难。'' This is
a feature, not a bug. \(\checkmark\) - PR-AUC: Same amplification
applies because PR-AUC =
\(\int_0^1 \text{Precision}(\text{Recall}^{-1}(t)) dt\), which at each
threshold involves conditional probabilities
\(\mathbb{P}(Z=1 \mid \hat{z}=1)\), requiring division by
\(\mathbb{P}(\hat{z}=1)\) (bounded below by factors involving \(\eta\)).
The document's claim that the same TV bound applies is correct at the
level of the joint distribution, though a rigorous PR-AUC bound would
require a more detailed argument. We flag this as a minor gap (see
Section 8). \(\checkmark\)

\textbf{Symmetric lower bound}: The theorem also has:
\[AUC(h_{\text{SCX}}) \geq AUC_{\text{base}} - \sqrt{\frac{\delta}{2}} \cdot \left(\frac{1}{\eta} + \frac{1}{1-\eta}\right)\]
This follows from the same TV bound applied in the opposite direction.
It means SCX cannot be much worse than baseline either --- the bound is
two-sided. \(\checkmark\)

\begin{center}\rule{0.5\linewidth}{0.5pt}\end{center}

\subsection{3. Theorem 3: Noise-Difficulty
Unidentifiability}\label{theorem-3-noise-difficulty-unidentifiability}

\subsubsection{3.1 Polished Statement}\label{polished-statement-2}

\textbf{Theorem 3 (Noise-Difficulty Unidentifiability).} For any
\(K_{\mathcal{Y}} \geq 2\) classification problem, any \(M \geq 1\)
experts, and any finite state space \(\mathcal{S}\), there exist two
data-generating processes \(\mathcal{P}_{\text{noise}}\) and
\(\mathcal{P}_{\text{hard}}\) such that:

\textbf{(i)} Under \(\mathcal{P}_{\text{noise}}\), label errors are
caused by uniform label noise (flips at rate \(\eta\)). \textbf{(ii)}
Under \(\mathcal{P}_{\text{hard}}\), all observed labels equal the true
labels, but some samples are intrinsically difficult (expert errors
arise from label ambiguity, not noise). \textbf{(iii)} The two processes
produce identical observable joint distributions:
\[\mathcal{P}_{\text{noise}}(x, y, \{f_m(x)\}) = \mathcal{P}_{\text{hard}}(x, y, \{f_m(x)\}), \quad \forall (x, y, \{f_m\})\]
\textbf{(iv)} Consequently, any algorithm \(\mathcal{A}\) that maps
\(n\) observations to per-sample noise flags has, for at least one of
the two worlds, an expected error rate \(\geq \eta\rho/2\) on the
ambiguous subset, where \(\eta\) is the noise rate and \(\rho\) is the
proportion of the ambiguous state.

\subsubsection{3.2 Construction (K=2) and
Verification}\label{construction-k2-and-verification}

\textbf{World A (Noise)}: State \(s_1\) (proportion \(\rho\)):
\(y^* \equiv 0\), label flipped to \(1\) with probability \(\eta\).
Expert accuracy on clean samples: \(1-\varepsilon_1\). State \(s_2\)
(proportion \(1-\rho\)): no noise, expert accuracy \(1-\varepsilon_2\).

\textbf{World B (Hard)}: State \(s_1\): \(y^* = 0\) with prob
\(1-\eta\), \(y^* = 1\) with prob \(\eta\). No label noise
(\(y = y^*\)). Experts biased to 0:
\(\mathbb{P}(f_m=0 \mid y^*=0) = 1-\varepsilon_1\),
\(\mathbb{P}(f_m=0 \mid y^*=1) = 1-\varepsilon_1\). State \(s_2\): same
as World A.

\textbf{Check}: \(\mathcal{P}(y=0 \mid s_1) = 1-\eta\) in both worlds.
\(\mathcal{P}(f_m=0 \mid s_1) = 1-\varepsilon_1\) in both worlds. Under
World A, \(y \perp f_m \mid s_1\) (noise independent of experts per A4).
Under World B, \(f_m \perp y^* \mid s_1\) by construction (expert 0-bias
independent of true label), hence \(f_m \perp y \mid s_1\) since
\(y=y^*\). Thus joint distributions match. \(\checkmark\)

\subsubsection{3.3 K \textgreater{} 2 Construction: Random
Experts}\label{k-2-construction-random-experts}

\textbf{FLAG (from verification, now FIXED)}: The original
K\textgreater2 construction used experts with non-trivial accuracy
conditional on true label, which caused \(\mathcal{P}(f_m \mid s_1)\) to
differ between worlds. The 2026-06-27 correction uses \textbf{fully
random experts} in World B.

\textbf{Corrected construction for \(K_{\mathcal{Y}} > 2\):}

\textbf{World A (Noise)}: State \(s_1\): \(y^* \equiv 0\), noise label
uniform over \(\{1,\dots,K_{\mathcal{Y}}-1\}\) with prob \(\eta\).
Expert: \(\mathbb{P}(f_m=0 \mid s_1) = 1-\varepsilon_1\),
\(\mathbb{P}(f_m=c \mid s_1) = \varepsilon_1/(K_{\mathcal{Y}}-1)\) for
\(c \neq 0\).

\[\mathcal{P}_{\text{noise}}(y=0 \mid s_1) = 1-\eta, \quad \mathcal{P}_{\text{noise}}(y=c \mid s_1) = \frac{\eta}{K_{\mathcal{Y}}-1}\]
\[\mathcal{P}_{\text{noise}}(f_m=0 \mid s_1) = 1-\varepsilon_1, \quad \mathcal{P}_{\text{noise}}(f_m=c \mid s_1) = \frac{\varepsilon_1}{K_{\mathcal{Y}}-1}\]

\textbf{World B (Hard)}: State \(s_1\): true label distribution
\(\mathbb{P}(y^*=0 \mid s_1) = 1-\eta\),
\(\mathbb{P}(y^*=c \mid s_1) = \eta/(K_{\mathcal{Y}}-1)\). Experts
\textbf{are random}: \(f_m \perp y^*\), with same marginal as World A.

The key difference from K=2: experts in World B are \textbf{fully
random} (do not depend on \(y^*\)), so \(f_m \perp y \mid s_1\) holds
automatically. This makes the joint distribution factorize identically
in both worlds. \(\checkmark\)

\subsubsection{3.4 Error Lower Bound
Derivation}\label{error-lower-bound-derivation}

The ambiguous subset is \(\{x \in s_1, y \neq 0\}\), of proportion
\(\rho \cdot \eta\). In World A, these are genuine noise samples (should
be flagged). In World B, these are genuine hard samples (should not be
flagged). Any algorithm flags fraction \(a\) of this subset. Then: -
Error in World A: \(\rho\eta(1-a)\) (false negatives) - Error in World
B: \(\rho\eta a\) (false positives) -
\(\max(\text{Error}_A, \text{Error}_B) \geq (\text{Error}_A + \text{Error}_B)/2 = \rho\eta/2\)

The factor 1/2 comes from the two-world average. This is a standard
minimax lower bound argument. When \(a = 1/2\) (random guessing on the
ambiguous set), both errors equal \(\rho\eta/2\), achieving the bound.
\(\checkmark\)

\textbf{Interpretation}: For \(\eta = 0.1\), \(\rho = 0.5\), the lower
bound is \(0.025\) --- small but strictly positive. The theorem
establishes qualitative unidentifiability (perfect distinction is
impossible), not a quantitative hardness threshold. \(\checkmark\)

\begin{center}\rule{0.5\linewidth}{0.5pt}\end{center}

\subsection{4. Theorem 4': Exact Constant Minimax
Optimality}\label{theorem-4-exact-constant-minimax-optimality}

\subsubsection{4.1 Polished Statement}\label{polished-statement-3}

\textbf{Theorem 4' (Exact Constant Minimax Optimality of SCX Noise
Detection).} Let assumptions (A1)-(A6) hold. Consider a single state
\(s\) with clean error rate \(p_0 = \mu_s\) and noise error rate
\(p_1 = 1 - C_{\text{bal}} \cdot \mu_s/(K_{\mathcal{Y}}-1)\), where
\(0 < p_0 < p_1 < 1\). Define:

\begin{itemize}
\tightlist
\item
  \(\kappa = C(\text{Bern}(p_0), \text{Bern}(p_1))\): Chernoff
  information
\item
  \(\theta^*\): unique solution of
  \(\text{KL}(\theta\|p_0) = \text{KL}(\theta\|p_1)\) (Chernoff point)
  \[\theta^* = \frac{\log\frac{1-p_0}{1-p_1}}{\log\frac{p_1(1-p_0)}{p_0(1-p_1)}}\]
\item
  \(\lambda_0^* = \log\frac{\theta^*(1-p_0)}{p_0(1-\theta^*)} > 0\),
  \(\lambda_1^* = \log\frac{\theta^*(1-p_1)}{p_1(1-\theta^*)} < 0\)
\item
  \(D^* = \lambda_0^* + |\lambda_1^*| = \log\frac{p_1(1-p_0)}{p_0(1-p_1)}\),
  \(s = |\lambda_1^*| / D^* \in (0,1)\)
\end{itemize}

Then:

\textbf{(a) SCX Achievability with Adaptive Threshold.} The SCX noise
detector using the adaptive threshold
\[\theta^\dagger = \theta^* + \frac{1}{M}\frac{\log((1-\eta)/\eta)}{D^*} + O(1/M^2)\]
achieves:
\[\boxed{\;\lim_{M\to\infty} e^{M\kappa} \cdot \sqrt{2\pi M} \cdot \bigl(1 - \text{F1}_{\text{SCX}}(\theta^\dagger)\bigr) = \frac{C_{\min}}{\eta}\;}\]
where the canonical minimax constant is:
\[\boxed{\;C_{\min} = \frac{\eta}{2} \left(\frac{1-\eta}{\eta}\right)^{s} \cdot \frac{1/\lambda_0^* + 1/|\lambda_1^*|}{\sqrt{\theta^*(1-\theta^*)}}\;}\]

\textbf{(b) Minimax Lower Bound.} For any noise detection algorithm
\(\mathcal{A}\):
\[\boxed{\;\liminf_{M\to\infty} e^{M\kappa} \cdot \sqrt{2\pi M} \cdot \bigl(1 - \text{F1}_{\mathcal{A}}\bigr) \geq \frac{C_{\min}}{\eta}\;}\]

\textbf{(c) Constant Optimality.} Since SCX with \(\theta^\dagger\)
achieves the limit with constant \(C_{\min}\), it is \textbf{exact
constant minimax optimal}.

\textbf{(d) Multi-state generalization.} For \(S\) states with
proportions \(\rho_s\), Chernoff information \(\kappa_s\), and per-state
constants \(C_s\):
\[\kappa_{\text{global}} = \min_s \kappa_s, \qquad C_{\text{global}} = \sum_{s: \kappa_s = \kappa_{\text{global}}} \rho_s C_s\]
and
\(\liminf e^{M\kappa_{\text{global}}}\sqrt{2\pi M}(1-\text{F1}_{\mathcal{A}}) \geq C_{\text{global}}/\eta\).

\subsubsection{4.2 Verification of
Constants}\label{verification-of-constants}

\textbf{CRITICAL: \(C_{\min}\) canonical form.} The verification report
identified three inconsistent \(C_{\min}\) formulas across the
manuscript files. \textbf{All inconsistencies have been resolved.} The
canonical \(C_{\min}\) is the Lemma E expression:

\[C_{\min} = \frac{\eta}{2} \left(\frac{1-\eta}{\eta}\right)^{s} \frac{1/\lambda_0^* + 1/|\lambda_1^*|}{\sqrt{\theta^*(1-\theta^*)}}\]

This is consistently used in: - Lemma E (eq. 45) \(\checkmark\) - Lemma
D, Theorem D.7 (explicitly references Lemma E) \(\checkmark\) - Theorem
4'(a) (this document) \(\checkmark\)

\textbf{Legacy formulas that were inconsistent (now deprecated)}: -
Architecture document Section 3.4 \(C_{\text{SCX}}\) formula (used
\(\max\) and lacked \(((1-\eta)/\eta)^s\) factor) -- REPLACED by Lemma E
form - Architecture document Section 4.3 \(C_{\min}\) formula (draft) --
REPLACED by Lemma E form

\textbf{Check: \(((1-\eta)/\eta)^s\) prefactor derivation.} - From Lemma
D.2:
\(\theta^\dagger = \theta^* + \frac{1}{M}\frac{\log((1-\eta)/\eta)}{D^*} + O(1/M^2)\).
- Taylor expand
\(M\cdot\text{KL}(\theta^\dagger\|p_0) = M\kappa + \lambda_0^*\frac{\log((1-\eta)/\eta)}{D^*} + O(1/M)\).
- Exponentiate:
\(\exp(-M\cdot\text{KL}(\theta^\dagger\|p_0)) = e^{-M\kappa} \cdot \left(\frac{1-\eta}{\eta}\right)^{-\lambda_0^*/D^*} \cdot (1+o(1))\).
- Similarly,
\(\exp(-M\cdot\text{KL}(\theta^\dagger\|p_1)) = e^{-M\kappa} \cdot \left(\frac{1-\eta}{\eta}\right)^{|\lambda_1^*|/D^*} \cdot (1+o(1))\).
- Since \(s = |\lambda_1^*|/D^*\) and \(\lambda_0^*/D^* = 1-s\), both
FPR and FNR contributions carry the \textbf{identical} factor
\(((1-\eta)/\eta)^s\), which factors out in the F1 expression.
\(\checkmark\)

\textbf{Check: Adaptive threshold constant matches lower bound.} - Lemma
D.7:
\(\lim e^{M\kappa}\sqrt{2\pi M}(1-\text{F1}_{\text{SCX}}(\theta^\dagger)) = \frac{((1-\eta)/\eta)^s}{2\sqrt{\theta^*(1-\theta^*)}}\left(\frac{1}{\lambda_0^*} + \frac{1}{|\lambda_1^*|}\right)\).
- Lemma E:
\(\liminf e^{M\kappa}\sqrt{2\pi M}(1-\text{F1}_{\mathcal{A}}) \geq \frac{C_{\min}}{\eta} = \frac{1}{2}\left(\frac{1-\eta}{\eta}\right)^s \frac{1/\lambda_0^* + 1/|\lambda_1^*|}{\sqrt{\theta^*(1-\theta^*)}}\).
- These are \textbf{identical expressions}. The factor of \(\eta\)
cancellation:
\(C_{\min}/\eta = \frac{\eta}{2\eta}(\frac{1-\eta}{\eta})^s \frac{1/\lambda_0^*+1/|\lambda_1^*|}{\sqrt{\theta^*(1-\theta^*)}} = \frac{1}{2}(\frac{1-\eta}{\eta})^s \frac{1/\lambda_0^*+1/|\lambda_1^*|}{\sqrt{\theta^*(1-\theta^*)}}\).
\(\checkmark\)

\subsubsection{4.3 Numerical Verification
Results}\label{numerical-verification-results}

Script \texttt{numerical\_verify.py} executed with results:

\begin{longtable}[]{@{}
  >{\raggedright\arraybackslash}p{(\linewidth - 18\tabcolsep) * \real{0.0462}}
  >{\raggedright\arraybackslash}p{(\linewidth - 18\tabcolsep) * \real{0.0538}}
  >{\raggedright\arraybackslash}p{(\linewidth - 18\tabcolsep) * \real{0.0538}}
  >{\raggedright\arraybackslash}p{(\linewidth - 18\tabcolsep) * \real{0.0615}}
  >{\raggedright\arraybackslash}p{(\linewidth - 18\tabcolsep) * \real{0.0769}}
  >{\raggedright\arraybackslash}p{(\linewidth - 18\tabcolsep) * \real{0.1615}}
  >{\raggedright\arraybackslash}p{(\linewidth - 18\tabcolsep) * \real{0.0923}}
  >{\raggedright\arraybackslash}p{(\linewidth - 18\tabcolsep) * \real{0.1308}}
  >{\raggedright\arraybackslash}p{(\linewidth - 18\tabcolsep) * \real{0.2000}}
  >{\raggedright\arraybackslash}p{(\linewidth - 18\tabcolsep) * \real{0.1231}}@{}}
\toprule\noalign{}
\begin{minipage}[b]{\linewidth}\raggedright
Case
\end{minipage} & \begin{minipage}[b]{\linewidth}\raggedright
\(p_0\)
\end{minipage} & \begin{minipage}[b]{\linewidth}\raggedright
\(p_1\)
\end{minipage} & \begin{minipage}[b]{\linewidth}\raggedright
\(\eta\)
\end{minipage} & \begin{minipage}[b]{\linewidth}\raggedright
\(\kappa\)
\end{minipage} & \begin{minipage}[b]{\linewidth}\raggedright
\(2\Delta^2\)/\(\kappa\)
\end{minipage} & \begin{minipage}[b]{\linewidth}\raggedright
\(C_{\min}\)
\end{minipage} & \begin{minipage}[b]{\linewidth}\raggedright
\(C_{\min}/\eta\)
\end{minipage} & \begin{minipage}[b]{\linewidth}\raggedright
Non-adapt/\(C_{\min}/\eta\)
\end{minipage} & \begin{minipage}[b]{\linewidth}\raggedright
Adapt matches?
\end{minipage} \\
\midrule\noalign{}
\endhead
\bottomrule\noalign{}
\endlastfoot
1 & 0.10 & 0.60 & 0.10 & 0.1696 & 2.95 & 0.4591 & 4.5915 & 1.70 & YES \\
2 & 0.20 & 0.50 & 0.30 & 0.0528 & 3.41 & 1.3768 & 4.5894 & 1.09 & YES \\
3 & 0.05 & 0.80 & 0.05 & 0.4574 & 2.46 & 0.1843 & 3.6864 & 2.41 & YES \\
4 & 0.10 & 0.60 & 0.50 & 0.1696 & 2.95 & 0.8348 & 1.6697 & 1.00 & YES \\
5 & 0.10 & 0.60 & 0.90 & 0.1696 & 2.95 & 0.5465 & 0.6072 & 1.62 & YES \\
\end{longtable}

\textbf{Key findings}: 1. Adaptive limit = \(C_{\min}/\eta\) to machine
precision (diff \(< 10^{-15}\)) for all cases. \(\checkmark\) 2.
Non-adaptive (naive \(\theta^*\)) is suboptimal by factors 1.00-2.41.
\(\checkmark\) 3. \(\kappa < 2\Delta^2\) for all cases (ratio
2.46-3.41), confirming the Chernoff rate is \textbf{slower} than the
Hoeffding bound. \(\checkmark\) 4. For \(\eta=0.5\), adaptive and
non-adaptive thresholds coincide (\(\theta^\dagger = \theta^*\)).
\(\checkmark\)

\subsubsection{4.4 Key Insight: The O(1/M) Threshold
Shift}\label{key-insight-the-o1m-threshold-shift}

The \(O(1/M)\) difference between \(\theta^*\) (Chernoff point) and
\(\theta^\dagger\) (adaptive threshold) produces an \textbf{O(1)
multiplicative factor} in the error probability via
\(\exp(-M\cdot\text{KL})\), not an \(o(1)\) correction. This factor
\(((1-\eta)/\eta)^s\) is essential for achieving the minimax lower
bound. The naive threshold \(\theta^*\) (which ignores \(\eta\)) is
suboptimal by this factor, and the gap can be up to \(2.41\times\) (Case
3).

\textbf{Historical correction}: Lemma D in an earlier version
incorrectly claimed that the O(1/M) shift produces only \(o(1)\)
effects. This was corrected in the 2026-06-27 revision. The current
Lemma D correctly handles the shift. \(\checkmark\)

\begin{center}\rule{0.5\linewidth}{0.5pt}\end{center}

\subsection{5. Theorem 5: Fixed-K Cluster
Consistency}\label{theorem-5-fixed-k-cluster-consistency}

\subsubsection{5.1 Polished Statement}\label{polished-statement-4}

\textbf{Theorem 5 (Fixed-K State Discovery Consistency).} Let
\(K_{\mathcal{S}}\) be fixed. Suppose:

\begin{enumerate}
\def\labelenumi{\arabic{enumi}.}
\tightlist
\item
  \textbf{Generative model}: \(\phi(x) = \mu_{s(x)} + \varepsilon\),
  where \(\varepsilon\) is zero-mean sub-Gaussian with variance proxy
  \(\sigma^2\), independent of \(s(x)\).
\item
  \textbf{Strong separation}:
  \(\Delta_{\min} = \min_{i \neq j} \|\mu_i - \mu_j\|_2 > 0\) satisfies
  \(\Delta_{\min}^2 / (\sigma^2 d_\phi) \geq C_0\) for a universal
  constant \(C_0\).
\item
  \textbf{Asymptotics}: \(n \to \infty\),
  \(n_{\min} = \min_k n_k \to \infty\).
\item
  \textbf{Algorithm}: Lloyd's k-means with \(R = C_R \log n\) random
  initializations, returning the solution with smallest \(W_n\).
\end{enumerate}

Then the estimated partition \(\hat{\mathcal{C}}^{(n)}\) satisfies:

\[\boxed{\;P\!\left(\hat{\mathcal{C}}^{(n)} \neq \mathcal{C}^* \text{ up to permutation}\right) \leq K_{\mathcal{S}} \cdot \exp\!\left(-c_1 \cdot \frac{n_{\min} \Delta_{\min}^2}{\sigma^2 d_\phi}\right) + o(1)\;}\]

where \(c_1 > 0\) is a universal constant. The \(o(1)\) term absorbs
exponentially small bias from the population minimizer
(\(\varepsilon_{\text{pop}}\)), initialization failure probability
(\(n^{-c}\)), and irreducible noise-boundary points
(\(\|\varepsilon\| \geq 3\Delta_{\min}/8\)).

\subsubsection{5.2 Verification of Hostile Review
Issues}\label{verification-of-hostile-review-issues}

\begin{longtable}[]{@{}
  >{\raggedright\arraybackslash}p{(\linewidth - 4\tabcolsep) * \real{0.2593}}
  >{\raggedright\arraybackslash}p{(\linewidth - 4\tabcolsep) * \real{0.2963}}
  >{\raggedright\arraybackslash}p{(\linewidth - 4\tabcolsep) * \real{0.4444}}@{}}
\toprule\noalign{}
\begin{minipage}[b]{\linewidth}\raggedright
Issue
\end{minipage} & \begin{minipage}[b]{\linewidth}\raggedright
Status
\end{minipage} & \begin{minipage}[b]{\linewidth}\raggedright
Resolution
\end{minipage} \\
\midrule\noalign{}
\endhead
\bottomrule\noalign{}
\endlastfoot
\textbf{Issue 1}: Reversed inequality in misclassification bound & FIXED
& Lemma 3 uses correct triangle inequality:
\(\|\phi - \theta_{j_k}\| \leq \Delta_{\min}/2 \leq \|\phi - \theta_{j'}\|\).
Both directions verified. \\
\textbf{Issue 2}: Positive exponent in uniform convergence & FIXED &
Localized empirical process with peeling gives prefactor 2 (no
\(n^{Kd_\phi}\) factor). Exponent \(-c_2 n t^2/(\sigma^2 d_\phi)\) is
explicitly negative for all \(n,t>0\). \\
\textbf{Issue 3}: NP-hard gap of k-means & FIXED & Lemma 4 proves
Lloyd's with \(R=\Omega(\log n)\) finds global minimizer under strong
separation. Honest gap discussion retained. \\
\textbf{Issue 4}: Lemma 5's incorrect lower bound & FIXED & Replaced
direct \(W(\mu)-W(\mu^*)\) computation with self-consistency argument
(Banach fixed-point). \\
\textbf{Issue 5}: Fixed vs.~scaling \(\Delta_{\min}\) & FIXED &
\(\Delta_{\min}\) is explicitly fixed, does not scale with \(n\). \\
\textbf{Issue 6c}: Triangle inequality noise handling & FIXED & Lemma 3
uses \(\|\varepsilon\| < 3\Delta_{\min}/8\) matching \(\Delta_{\min}/2\)
bounds. \\
\textbf{Issue 7}: Covering number vs VC dimension & FIXED & Uses
covering number directly (Lemma S3). \\
\textbf{Issue 10}: Chi-squared bound for sub-Gaussian & FIXED & Lemma S1
uses correct sub-Gaussian norm bound (Vershynin 2018, Theorem 3.1.1). \\
\end{longtable}

\textbf{Explicit exponent negativity check}: The exponent
\(-c_1 \cdot n_{\min} \Delta_{\min}^2/(\sigma^2 d_\phi)\) is always
negative for \(n_{\min} > 0\), \(\Delta_{\min} > 0\),
\(\sigma^2 < \infty\), \(d_\phi < \infty\). The constant
\(c_1 = c_2 K_{\mathcal{S}}/64\) where
\(c_2 = c_6 \lambda^2/(128 C_L^2)\), \(\lambda = \pi_{\min}/2 > 0\). All
constants are positive. \(\checkmark\)

\subsubsection{5.3 NP-Hard Gap Discussion}\label{np-hard-gap-discussion}

\textbf{Context}: k-means is NP-hard in the worst case (Aloise et al.,
2009). Lemma 2 assumes access to the global empirical minimizer
\(\hat{\theta}_n\).

\textbf{Resolution under strong separation}: Lemma 4 proves that under
the strong separation condition
(\(\Delta_{\min}^2/(\sigma^2 d_\phi) \geq C_0\)), the k-means landscape
has a unique local minimum within \(\Delta_{\min}/4\) of the true
centers. Lloyd's algorithm with \(R = C_R \log n\) random
initializations contracts toward the global minimizer with probability
\(1 - n^{-c}\).

\textbf{Without strong separation}: The guarantee degrades. In this
regime, the theorem only claims existence of some polynomial-time
algorithm finding an \(\varepsilon\)-approximate solution (e.g., via
Kumar \& Kannan, 2010 or Ostrovsky et al., 2013). This gap is honestly
stated and does not affect the theorem's validity under its assumptions.

\begin{center}\rule{0.5\linewidth}{0.5pt}\end{center}

\subsection{6. Proposition 6: Bootstrap Stability Diagnostic for Feature
Strength}\label{proposition-6-bootstrap-stability-diagnostic-for-feature-strength}

\subsubsection{6.1 Statement}\label{statement}

\textbf{Proposition 6 (Clustering Stability Criterion for SCX Feature
Strength).} For feature vectors \(\{\phi(x_i)\}_{i=1}^N\) and
\(K_{\mathcal{S}}\) states, define clustering stability
\(S(\Phi, K_{\mathcal{S}})\) as the mean adjusted Rand index (ARI)
between k-means clusterings on the full data and on bootstrap resamples.
Then:

\textbf{(a) Strong features \(\implies\) high stability.} If
\(\Delta_{\min}^2/\sigma^2 > C_1(d + \log N)/n_{\min}\), then with
probability \(\geq 1 - O(N^{-1})\):
\[S(\Phi, K_{\mathcal{S}}) > 1 - \varepsilon, \quad \varepsilon = O(e^{-c \cdot n_{\min} \Delta_{\min}^2/\sigma^2})\]

\textbf{(b) Weak/absent state structure \(\implies\) low stability.} If
all \(\mu_k\) are equal (no state structure), then:
\[\mathbb{E}[S(\Phi, K_{\mathcal{S}})] = O(K_{\mathcal{S}}/\sqrt{N}) \quad \text{(near 0)}\]

\textbf{(c) Diagnostic for SCX reliability.} If
\(S(\Phi, K_{\mathcal{S}}) < \tau\) (with \(\tau = 0.7\) as recommended
heuristic), then with high probability the features are too weak for SCX
to significantly outperform the loss baseline. Conversely, if
\(S(\Phi, K_{\mathcal{S}}) > \tau\), features are not the bottleneck
(though other factors like expert redundancy still matter).

\subsubsection{6.2 Connection to Theorem
2}\label{connection-to-theorem-2}

The stability diagnostic operationalizes Theorem 2's
\(\delta = I(\phi; S)\):

\begin{itemize}
\tightlist
\item
  \textbf{Low stability} (\(S < 0.5\)): The estimated state partition
  \(\hat{S}\) from k-means is essentially random. The mutual information
  \(I(\phi; \hat{S})\) is near zero, implying \(\delta \approx 0\). By
  Theorem 2,
  \(\text{F1}_{\text{SCX}} \leq \text{F1}_{\text{base}} + o(1)\) --- SCX
  cannot improve over loss baseline.
\item
  \textbf{High stability} (\(S > 0.7\)): k-means produces a reproducible
  partition. The effective mutual information \(I(\phi; \hat{S})\) is
  bounded below by a function of \(S\), so Theorem 2's bound is not
  tight --- SCX may improve over the baseline.
\item
  \textbf{Intermediate stability} (\(0.5 \leq S \leq 0.7\)):
  Transitional regime. Theorem 2 provides a non-trivial but weak bound.
\end{itemize}

\textbf{Caveat}: The connection is empirical, not formally proven.
Proposition 6's Part (c) provides a heuristic threshold, not a theorem.
Unlike the BBP spectral proxy (which attempted a formal connection and
failed), the stability diagnostic directly tests what SCX does (k-means
clustering) without distributional assumptions. See Section 8 for
limitations.

\begin{center}\rule{0.5\linewidth}{0.5pt}\end{center}

\subsection{7. Cross-Theorem Consistency
Audit}\label{cross-theorem-consistency-audit}

\subsubsection{7.1 Notation Conflict
Resolution}\label{notation-conflict-resolution}

\begin{longtable}[]{@{}
  >{\raggedright\arraybackslash}p{(\linewidth - 14\tabcolsep) * \real{0.1250}}
  >{\raggedright\arraybackslash}p{(\linewidth - 14\tabcolsep) * \real{0.1094}}
  >{\raggedright\arraybackslash}p{(\linewidth - 14\tabcolsep) * \real{0.1094}}
  >{\raggedright\arraybackslash}p{(\linewidth - 14\tabcolsep) * \real{0.1094}}
  >{\raggedright\arraybackslash}p{(\linewidth - 14\tabcolsep) * \real{0.1250}}
  >{\raggedright\arraybackslash}p{(\linewidth - 14\tabcolsep) * \real{0.1094}}
  >{\raggedright\arraybackslash}p{(\linewidth - 14\tabcolsep) * \real{0.1250}}
  >{\raggedright\arraybackslash}p{(\linewidth - 14\tabcolsep) * \real{0.1875}}@{}}
\toprule\noalign{}
\begin{minipage}[b]{\linewidth}\raggedright
Symbol
\end{minipage} & \begin{minipage}[b]{\linewidth}\raggedright
Thm 1
\end{minipage} & \begin{minipage}[b]{\linewidth}\raggedright
Thm 2
\end{minipage} & \begin{minipage}[b]{\linewidth}\raggedright
Thm 3
\end{minipage} & \begin{minipage}[b]{\linewidth}\raggedright
Thm 4'
\end{minipage} & \begin{minipage}[b]{\linewidth}\raggedright
Thm 5
\end{minipage} & \begin{minipage}[b]{\linewidth}\raggedright
Prop 6
\end{minipage} & \begin{minipage}[b]{\linewidth}\raggedright
Resolution
\end{minipage} \\
\midrule\noalign{}
\endhead
\bottomrule\noalign{}
\endlastfoot
\(K\) & \(K_{\mathcal{Y}}\) classes & \(K_{\mathcal{S}}\) states &
\(K_{\mathcal{Y}}\) classes & \(K_{\mathcal{Y}}\) classes &
\(K_{\mathcal{S}}\) states & \(K\) clusters & Disambiguated as
\(K_{\mathcal{Y}}, K_{\mathcal{S}}\) in this doc \\
\(\theta\) & Detection threshold & --- & --- & Threshold/Chernoff point
& --- & --- & No conflict; same meaning \\
\(\mu\) & Clean error bound & --- & --- & Clean error rate \(p_0\) &
Cluster centers & --- & Different meanings; context disambiguated \\
\(\eta\) & Noise rate & Noise rate & Noise/ambiguity rate & Noise rate &
--- & --- & Same meaning in all \\
\(\Delta\) & Separation gap & --- & --- & --- & Min center separation &
Min center separation & Different (scalar vs gap); context
disambiguated \\
\(C(\cdot)\) & Consensus score & State consistency & --- & Chernoff info
& --- & --- & Disambiguated by argument \\
\end{longtable}

\subsubsection{7.2 Assumption Consistency}\label{assumption-consistency}

\begin{longtable}[]{@{}
  >{\raggedright\arraybackslash}p{(\linewidth - 12\tabcolsep) * \real{0.2000}}
  >{\raggedright\arraybackslash}p{(\linewidth - 12\tabcolsep) * \real{0.1273}}
  >{\raggedright\arraybackslash}p{(\linewidth - 12\tabcolsep) * \real{0.1273}}
  >{\raggedright\arraybackslash}p{(\linewidth - 12\tabcolsep) * \real{0.1273}}
  >{\raggedright\arraybackslash}p{(\linewidth - 12\tabcolsep) * \real{0.1455}}
  >{\raggedright\arraybackslash}p{(\linewidth - 12\tabcolsep) * \real{0.1273}}
  >{\raggedright\arraybackslash}p{(\linewidth - 12\tabcolsep) * \real{0.1455}}@{}}
\toprule\noalign{}
\begin{minipage}[b]{\linewidth}\raggedright
Assumption
\end{minipage} & \begin{minipage}[b]{\linewidth}\raggedright
Thm 1
\end{minipage} & \begin{minipage}[b]{\linewidth}\raggedright
Thm 2
\end{minipage} & \begin{minipage}[b]{\linewidth}\raggedright
Thm 3
\end{minipage} & \begin{minipage}[b]{\linewidth}\raggedright
Thm 4'
\end{minipage} & \begin{minipage}[b]{\linewidth}\raggedright
Thm 5
\end{minipage} & \begin{minipage}[b]{\linewidth}\raggedright
Prop 6
\end{minipage} \\
\midrule\noalign{}
\endhead
\bottomrule\noalign{}
\endlastfoot
A1: Disjoint training & REQUIRED & Implicit & Used in construction &
REQUIRED & N/A & N/A \\
A2: Conditional independence & REQUIRED & Implicit & Used in
construction & REQUIRED & N/A & N/A \\
A3: Bounded loss & REQUIRED & N/A & N/A & REQUIRED & N/A & N/A \\
A4: Uniform noise & REQUIRED & Implicit & REQUIRED for K\textgreater2 &
REQUIRED & N/A & N/A \\
A5: State homogeneity & REQUIRED & N/A & Sufficient for break & REQUIRED
& Used (i.i.d. within state) & N/A \\
A6: Balanced errors & REQUIRED & N/A & Sufficient for break & REQUIRED
(defines \(p_1\)) & N/A & N/A \\
\end{longtable}

\textbf{No conflicts}: Assumptions are consistently used across
theorems. Theorem 2 does not require A1-A6 explicitly but assumes the
SCX pipeline is applied; its negative result holds regardless of whether
A1-A6 are satisfied. Theorem 5 and Proposition 6 operate on features
only and do not require the label-noise assumptions. \(\checkmark\)

\subsubsection{7.3 Numerical Consistency: Three Parameter
Sets}\label{numerical-consistency-three-parameter-sets}

We verify that Theorem 1's bound \(\leq\) Theorem 4'\,'s exact
asymptotics \(\leq\) empirical results (where available).

\textbf{Parameter Set 1}: \(K_{\mathcal{Y}}=10\), \(\mu_s=0.2\),
\(\eta=0.1\), \(M=20\), \(C_{\text{bal}}=1\) (from Theorem 1 Tightness
section). - Thm 1 bound:
\(\text{F1} \geq 1 - 10 \times \exp(-2 \times 20 \times 0.389^2) = 1 - 10 \times e^{-6.05} > 0.976\).
- Thm 4' rate: \(\exp(-M\kappa)/\sqrt{M}\) where
\(\kappa = \text{KL}(\theta^*\|p_0) \approx 0.112\) (Chernoff info when
\(p_0=0.2\), \(p_1=0.911\)), giving
\(e^{-20 \times 0.112} = e^{-2.24} \approx 0.106\), slower than Thm 1's
\(e^{-6.05}\). The tighter bound from Theorem 4' is actually
\textbf{weaker} for finite \(M\) because \(\kappa < 2\Delta^2\). This is
expected: Theorem 1 provides a finite-sample Hoeffding guarantee, while
Theorem 4' provides the exact asymptotic constant that is tighter in the
\(M \to \infty\) limit but looser for moderate \(M\). -
\textbf{Consistency}: Both bounds are valid (Thm 1 is an inequality, Thm
4' is an asymptotic equality). No contradiction. \(\checkmark\)

\textbf{Parameter Set 2}: \(K_{\mathcal{Y}}=2\), \(\mu_s=\varepsilon\),
\(\theta=0.5\) (symmetric experts). - Thm 1:
\(\text{F1} \geq 1 - \frac{1}{\eta}\exp(-2M(1/2 - \varepsilon)^2)\). -
Thm 4': For K=2, \(p_1 = 1 - \mu_s = 1-\varepsilon\) (since
\(C_{\text{bal}}=1\), \(K_{\mathcal{Y}}-1=1\)). The Chernoff information
\(\kappa = C(\text{Bern}(\varepsilon), \text{Bern}(1-\varepsilon)) = \log 2 - H(\varepsilon)\)
where \(H\) is binary entropy. For \(\varepsilon=0.2\),
\(\kappa \approx 0.278\), \(2(0.5-0.2)^2 = 0.18\). So Thm 4'\,'s
\(\kappa >\) Thm 1's Hoeffding exponent. This is consistent: for K=2,
the Chernoff rate can be faster than the Hoeffding bound. \(\checkmark\)

\textbf{Parameter Set 3}: CIFAR-10 experiment (\(K_{\mathcal{Y}}=10\),
\(\mu_s\approx 0.45\), \(\eta=0.1\), \(M=20\)). - Thm 1 bound:
\(\text{F1} \geq 1 - \frac{1}{0.1}\exp(-2\times20\times0.25^2) = 1 - 10\times e^{-2.5} \approx 0.18\).
- Thm 4' exact constant: \(p_0=0.45\), \(p_1=1-0.45/9=0.95\),
\(\kappa = \text{KL}(\theta^*\|0.45) \approx 0.496\). For \(M=20\),
\(e^{-20\times0.496} = e^{-9.92} \approx 4.9\times10^{-5}\). - Empirical
F1: 0.617. Both bounds are satisfied (0.617 \textgreater{} 0.18 from Thm
1, and Thm 4' asymptotics would apply for larger \(M\)). \(\checkmark\)

\subsubsection{7.4 Sign Error
Verification}\label{sign-error-verification}

Every inequality direction across all theorems has been checked:

\begin{longtable}[]{@{}
  >{\raggedright\arraybackslash}p{(\linewidth - 6\tabcolsep) * \real{0.2381}}
  >{\raggedright\arraybackslash}p{(\linewidth - 6\tabcolsep) * \real{0.2619}}
  >{\raggedright\arraybackslash}p{(\linewidth - 6\tabcolsep) * \real{0.2619}}
  >{\raggedright\arraybackslash}p{(\linewidth - 6\tabcolsep) * \real{0.2381}}@{}}
\toprule\noalign{}
\begin{minipage}[b]{\linewidth}\raggedright
Location
\end{minipage} & \begin{minipage}[b]{\linewidth}\raggedright
Inequality
\end{minipage} & \begin{minipage}[b]{\linewidth}\raggedright
Direction
\end{minipage} & \begin{minipage}[b]{\linewidth}\raggedright
Verified
\end{minipage} \\
\midrule\noalign{}
\endhead
\bottomrule\noalign{}
\endlastfoot
Thm 1, Lemma 1 & \(\mathbb{E}[C \mid \text{clean}] \leq \mu_s\) &
\(\leq\) (upper bound) & \(\checkmark\) \\
Thm 1, Lemma 1 &
\(\mathbb{E}[C \mid \text{noise}] \geq 1 - \mu_s/(K_{\mathcal{Y}}-1)\) &
\(\geq\) (lower bound) & \(\checkmark\) \\
Thm 1, Lemma 2 &
\(\mathbb{P}(C > \theta \mid \text{clean}) \leq \exp(-2M(\theta-\mu_s)^2)\)
& \(\leq\) (tail bound) & \(\checkmark\) \\
Thm 1, Lemma 3 &
\(\mathbb{P}(C > \theta \mid \text{noise}) \geq 1 - \exp(-2M(1 - C_{\text{bal}}\mu_s/(K_{\mathcal{Y}}-1) - \theta)^2)\)
& \(\geq\) (1 - upper bound on error) & \(\checkmark\) \\
Thm 1, F1 bound &
\(\text{F1} \geq 1 - \delta_1 - \frac{1-\eta}{\eta}\delta_2\) & \(\geq\)
(lower bound) & \(\checkmark\) \\
Thm 1, Chernoff &
\(\mathbb{P}(C \leq \theta \mid \text{noise}) \leq \exp(-M \cdot \text{KL}(\theta \| p_1))\)
& \(\leq\) (lower tail) & \(\checkmark\) \\
Thm 2, AUC &
\(AUC(h_{\text{SCX}}) \leq AUC_{\text{base}} + \sqrt{\delta/2} \cdot (1/\eta + 1/(1-\eta))\)
& \(\leq\) (upper bound) & \(\checkmark\) \\
Thm 2, F1 &
\(F1(h_{\text{SCX}}) \leq F1_{\text{base}} + C_F\sqrt{\delta/2}\) &
\(\leq\) (upper bound) & \(\checkmark\) \\
Thm 5, Lemma 2 &
\(P(\|\hat{\theta}_n - \theta^*\| \geq t) \leq 2\exp(-c_2 n t^2/(\sigma^2 d_\phi))\)
& \(\leq\) (probability bound) & \(\checkmark\) \\
Thm 5, Lemma 3 &
\(\|\phi - \theta_{j_k}\| \leq \Delta_{\min}/2 \leq \|\phi - \theta_{j'}\|\)
& \(\leq\) on both sides & \(\checkmark\) \\
Thm 4', lower bound &
\(\liminf e^{M\kappa}\sqrt{2\pi M}(1-\text{F1}_{\mathcal{A}}) \geq C_{\min}/\eta\)
& \(\geq\) (lower bound) & \(\checkmark\) \\
Thm 4', achievability &
\(\lim e^{M\kappa}\sqrt{2\pi M}(1-\text{F1}_{\text{SCX}}) = C_{\min}/\eta\)
& \(=\) (exact limit) & \(\checkmark\) \\
Thm 3, error &
\(\max(\text{Error}_{\text{noise}}, \text{Error}_{\text{hard}}) \geq \eta\rho/2\)
& \(\geq\) (lower bound) & \(\checkmark\) \\
\end{longtable}

\textbf{No sign errors found.} All inequality directions are
mathematically correct. \(\checkmark\)

\begin{center}\rule{0.5\linewidth}{0.5pt}\end{center}

\subsection{8. Known Gaps and
Limitations}\label{known-gaps-and-limitations}

\subsubsection{8.1 Proven vs.~Conjectured}\label{proven-vs.-conjectured}

\begin{longtable}[]{@{}
  >{\raggedright\arraybackslash}p{(\linewidth - 4\tabcolsep) * \real{0.3182}}
  >{\raggedright\arraybackslash}p{(\linewidth - 4\tabcolsep) * \real{0.3636}}
  >{\raggedright\arraybackslash}p{(\linewidth - 4\tabcolsep) * \real{0.3182}}@{}}
\toprule\noalign{}
\begin{minipage}[b]{\linewidth}\raggedright
Claim
\end{minipage} & \begin{minipage}[b]{\linewidth}\raggedright
Status
\end{minipage} & \begin{minipage}[b]{\linewidth}\raggedright
Notes
\end{minipage} \\
\midrule\noalign{}
\endhead
\bottomrule\noalign{}
\endlastfoot
Thm 1: F1 bound under A1-A6 & \textbf{Proven} & Complete proof with
lemmas \\
Thm 1: Chernoff form of bound & \textbf{Proven} & Standard Chernoff
bound \\
Thm 1: Optimal threshold formula & \textbf{Proven} & Derivation from
separation gap \\
Thm 2: AUC bound & \textbf{Proven} & Complete Pinsker-TV chain \\
Thm 2: PR-AUC bound & \textbf{Partially proven} & TV bound holds for
joint distribution; PR-AUC requires conditioning on decision threshold,
which adds complexity. The bound is valid at the level claimed, but a
fully rigorous PR-AUC bound would need additional steps. \\
Thm 2: F1 bound & \textbf{Proven under operating range} & Lipschitz
constant \(C_F\) is well-defined; its bound of \(\leq 2\) requires
precision, recall \(\geq 0.1\), which is the typical operating range but
not proven to hold universally. \\
Thm 3: K=2 construction & \textbf{Proven} & Complete \\
Thm 3: K\textgreater2 construction & \textbf{Proven} & Corrected
2026-06-27 \\
Thm 3: Error lower bound \(\eta\rho/2\) & \textbf{Proven} & Standard
minimax argument \\
Thm 4': Chernoff point closed form & \textbf{Proven} & Lemma C \\
Thm 4': Bahadur-Rao asymptotics & \textbf{Proven} & Lemma A, complete
derivation \\
Thm 4': F1 asymptotic expansion & \textbf{Proven} & Lemma B, with error
bounds \\
Thm 4': Adaptive threshold optimality & \textbf{Proven} & Lemmas D + E,
numerical verification \\
Thm 4': Multi-state aggregation & \textbf{Proven} & Lemma F \\
Thm 5: Population minimizer proximity & \textbf{Proven} & Lemma 1,
self-consistency argument \\
Thm 5: Empirical minimizer convergence & \textbf{Proven} & Lemma 2,
localized empirical process \\
Thm 5: Lloyd's success & \textbf{Proven} & Lemma 4 \\
Thm 5: Final bound & \textbf{Proven} & Combined from lemmas \\
Prop 6: Part (a) strong \(\to\) high stability & \textbf{Proven} & Via
k-means stability literature (Shamir \& Tishby, 2009) \\
Prop 6: Part (b) weak \(\to\) low stability & \textbf{Proven} & ARI
expectation for random partitions \\
Prop 6: Part (c) diagnostic & \textbf{Heuristic} & The \(\tau=0.7\)
threshold is empirically motivated, not derived from first principles \\
Prop 6: Connection to Thm 2 & \textbf{Empirical} & Not formally proven;
based on the intuition that stability \(\implies\) identifiable state
structure \(\implies\) Thm 2's \(\delta\) not near zero \\
\end{longtable}

\subsubsection{8.2 Necessary vs.~Sufficient
Assumptions}\label{necessary-vs.-sufficient-assumptions}

\begin{longtable}[]{@{}
  >{\raggedright\arraybackslash}p{(\linewidth - 6\tabcolsep) * \real{0.2444}}
  >{\raggedright\arraybackslash}p{(\linewidth - 6\tabcolsep) * \real{0.2444}}
  >{\raggedright\arraybackslash}p{(\linewidth - 6\tabcolsep) * \real{0.2889}}
  >{\raggedright\arraybackslash}p{(\linewidth - 6\tabcolsep) * \real{0.2222}}@{}}
\toprule\noalign{}
\begin{minipage}[b]{\linewidth}\raggedright
Assumption
\end{minipage} & \begin{minipage}[b]{\linewidth}\raggedright
Necessary?
\end{minipage} & \begin{minipage}[b]{\linewidth}\raggedright
Sufficient?
\end{minipage} & \begin{minipage}[b]{\linewidth}\raggedright
Evidence
\end{minipage} \\
\midrule\noalign{}
\endhead
\bottomrule\noalign{}
\endlastfoot
A1 (disjoint training) & For Thm 1's exponential rate & Yes, with A2-A6
& Without A1, experts may correlate, breaking Hoeffding \\
A2 (conditional independence) & For Thm 1's concentration & Follows from
A1 & --- \\
A3 (bounded loss) & For Hoeffding & Yes & Can be relaxed to
sub-Gaussian \\
A4 (uniform noise) & For Thm 1's Lemma 1 & Yes, with A1-A3, A5-A6 &
Without A4, noise-consensus relationship changes \\
A5 (state homogeneity) & For per-state analysis & Yes, with A1-A4, A6 &
Without A5, state-level guarantees fail \\
A6 (balanced errors) & For Thm 1's TPR bound & Yes, with A1-A5 & Without
A6, \(C_{\text{bal}}\) can be large \\
Strong separation (Thm 5) & For polynomial-time recovery & Yes & Without
it, NP-hard gap remains \\
Sub-Gaussian noise (Thm 5) & For exponential rates & Can be relaxed &
Bounded noise gives slower \(\exp(-c n)\) instead of \(\exp(-c n)\) but
same form \\
Fixed \(K\) (Thm 5) & For current proof & For fixed-\(K\) rates &
Growing \(K\) needs separate analysis \\
\end{longtable}

\subsubsection{8.3 Edge Cases Not Fully
Covered}\label{edge-cases-not-fully-covered}

\begin{enumerate}
\def\labelenumi{\arabic{enumi}.}
\item
  \textbf{\(\eta \to 0\) or \(\eta \to 1\) in Theorem 4'}: The
  asymptotic expansion
  \(1-\text{F1} \approx \text{FNR}/2 + (1-\eta)\text{FPR}/(2\eta)\)
  breaks down when \(\eta \ll \eta_{\min}(M)\) (defined as
  \(\frac{e^{-M\kappa_0}}{2\lambda_0^*\sqrt{2\pi M\theta(1-\theta)}}\)).
  For realistic \(M\), \(\eta_{\min}(M)\) is exponentially small, so
  this only matters for extreme sparsity. Lemma B Section B.6.1 provides
  explicit validity conditions.
\item
  \textbf{\(p_0 = 0\) (perfect experts) in Theorem 4'}: The KL
  divergence \(\text{KL}(\theta\|0) = \infty\) for any \(\theta>0\),
  making \(\kappa = \infty\). The lower bound becomes vacuous
  (\(0 \geq 0\)), which is correct --- perfect experts make noise
  detection trivial. This edge case is discussed in Lemma E.
\item
  \textbf{\(K_{\mathcal{Y}} = 2\) with \(C_{\text{bal}} > 1\) in Theorem
  1}: When \(K_{\mathcal{Y}} = 2\),
  \(p_1 = 1 - C_{\text{bal}} \cdot \mu_s/1 = 1 - C_{\text{bal}}\mu_s\).
  For \(C_{\text{bal}} > 1\), if \(\mu_s > 1/C_{\text{bal}}\), then
  \(p_1 < \mu_s\) and the detection gap vanishes. The theorem requires
  \(p_1 > p_0\), i.e., \(\mu_s < 1/(C_{\text{bal}}+1)\) for K=2. This is
  a non-trivial restriction not discussed in the source files.
\item
  \textbf{Theorem 2 PR-AUC bound rigor}: The PR-AUC bound follows the
  same argument as the AUC bound, noting that PR-AUC at each threshold
  equals \(\mathbb{P}(Z=1 \mid \hat{z}=1)\), which involves conditional
  probabilities. The TV bound for AUC uses
  \(TV(P(\cdot|Z=1), Q(\cdot|Z=1))\) for the product measure of two
  conditionally independent scores. For PR-AUC, the argument requires
  additional steps because PR-AUC is an integral over thresholds of
  precision (\(\mathbb{P}(Z=1 \mid \hat{z}=1)\)), not a single
  expectation. The source files treat this at the same level of rigor as
  the AUC case; a fully rigorous treatment would require a more detailed
  argument.
\item
  \textbf{Proposition 6 threshold \(\tau = 0.7\)}: This value is
  borrowed from the Cohen's kappa agreement literature (Landis \& Koch,
  1977), not derived from SCX-specific theory. The calibration per
  domain (Section 4.2 of the source file) addresses this, but the
  threshold remains heuristic.
\end{enumerate}

\begin{center}\rule{0.5\linewidth}{0.5pt}\end{center}

\subsection{9. SI Writing Guide}\label{si-writing-guide}

\subsubsection{9.1 Source File to SI Section
Mapping}\label{source-file-to-si-section-mapping}

\begin{longtable}[]{@{}
  >{\raggedright\arraybackslash}p{(\linewidth - 8\tabcolsep) * \real{0.1667}}
  >{\raggedright\arraybackslash}p{(\linewidth - 8\tabcolsep) * \real{0.1364}}
  >{\raggedright\arraybackslash}p{(\linewidth - 8\tabcolsep) * \real{0.2273}}
  >{\raggedright\arraybackslash}p{(\linewidth - 8\tabcolsep) * \real{0.2879}}
  >{\raggedright\arraybackslash}p{(\linewidth - 8\tabcolsep) * \real{0.1818}}@{}}
\toprule\noalign{}
\begin{minipage}[b]{\linewidth}\raggedright
SI Section
\end{minipage} & \begin{minipage}[b]{\linewidth}\raggedright
Content
\end{minipage} & \begin{minipage}[b]{\linewidth}\raggedright
Source File(s)
\end{minipage} & \begin{minipage}[b]{\linewidth}\raggedright
Translation Needs
\end{minipage} & \begin{minipage}[b]{\linewidth}\raggedright
Est. Pages
\end{minipage} \\
\midrule\noalign{}
\endhead
\bottomrule\noalign{}
\endlastfoot
A.1 Notation and Assumptions & Unified glossary, A1-A6, dependency graph
& All theorem files & Minor polish & 2 \\
A.2 Theorem 1 & Statement, Lemma 1-3, F1 derivation, Chernoff appendix &
\texttt{01\_noise\_detection\_guarantee.md} & Remove \(C_{\text{bal}}\)
generalization discussion (keep formal). & 6 \\
A.3 Theorem 2 & Statement, Pinsker-TV chain, AUC/F1 bounds &
\texttt{02\_weak\_feature\_failure.md} & Streamline DermaMNIST
experimental references. & 5 \\
A.4 Theorem 3 & Construction, equivalence proof, K\textgreater2
extension & \texttt{03\_unidentifiability\_theorem.md} & Compress
Dawid-Skene comparison. & 4 \\
A.5 Theorem 4' & Bahadur-Rao, Chernoff information, adaptive threshold,
Lemmas A-F & \texttt{exact\_constant\_minimax.md},
\texttt{lemma\_AB\_*.md}, \texttt{lemma\_CD\_*.md},
\texttt{lemma\_EF\_*.md} & Heavy compression needed. Lemmas A-F each in
separate SI subsections. & 12 \\
A.6 Theorem 5 & Cluster consistency, Lemmas 1-4 &
\texttt{cluster\_consistency\_v3.md} & Includes hostile review rebuttal
notes. Keep concise. & 8 \\
A.7 Proposition 6 & Stability diagnostic &
\texttt{feature\_strength\_via\_stability.md} & Compress BBP failure
analysis (Section 1). & 3 \\
A.8 Numerical verification & Table of constants, code &
\texttt{numerical\_verify.py}, \texttt{verification\_exact\_constant.md}
& Minimal --- just final table. & 1 \\
\end{longtable}

\subsubsection{9.2 Translation Notes}\label{translation-notes}

\textbf{For LaTeX translation}: - All inline math is already in
LaTeX-compatible notation. - Theorem environments should use
\texttt{\textbackslash{}begin\{theorem\}...\textbackslash{}end\{theorem\}}.
- Lemma B's F1 expansion should be presented as asymptotic series. -
Lemma E's Bayes test derivation should be the SI centerpiece for Theorem
4'. - The verification report's numerical table (Section 9 of this
document) should be included as a verification table.

\textbf{Direct inclusion possible}: - Notation tables (Section 0.1 of
this document). - Assumption catalog (Section 0.2). - Lemma statements
(not full proofs, just statements). - Numerical verification table.

\textbf{Needs LaTeX-only reformatting}: - Theorem proofs (convert
\boxed{} environments to align environments). - Long derivations (Lemma
A Stirling expansion, Lemma D adaptive threshold derivation).

\subsubsection{9.3 Cross-References}\label{cross-references}

\begin{itemize}
\tightlist
\item
  Theorem 1 references Theorem 3 for necessity justification.
\item
  Theorem 2 references Proposition 6 for practical diagnostic.
\item
  Theorem 4' references Lemma B from Theorem 1's proof chain.
\item
  Theorem 5 references Theorem 2 as the negative counterpart.
\item
  All theorems reference the assumption catalog (A1-A6).
\end{itemize}

\subsubsection{9.4 Estimated Total SI
Pages}\label{estimated-total-si-pages}

\begin{longtable}[]{@{}ll@{}}
\toprule\noalign{}
Section & Pages \\
\midrule\noalign{}
\endhead
\bottomrule\noalign{}
\endlastfoot
A.1 Notation and Assumptions & 2 \\
A.2 Theorem 1 & 6 \\
A.3 Theorem 2 & 5 \\
A.4 Theorem 3 & 4 \\
A.5 Theorem 4' (Lemmas A-F) & 12 \\
A.6 Theorem 5 & 8 \\
A.7 Proposition 6 & 3 \\
A.8 Numerical Verification & 1 \\
\textbf{Total} & \textbf{41} \\
\end{longtable}

\begin{center}\rule{0.5\linewidth}{0.5pt}\end{center}

\subsection{10. References}\label{references}

\begin{enumerate}
\def\labelenumi{\arabic{enumi}.}
\tightlist
\item
  Bahadur, R. R., \& Rao, R. R. (1960). On deviations of the sample
  mean. \emph{Annals of Mathematical Statistics}, 31(4), 1015-1027.
\item
  Chernoff, H. (1952). A measure of asymptotic efficiency for tests of a
  hypothesis based on the sum of observations. \emph{Annals of
  Mathematical Statistics}, 23(4), 493-507.
\item
  Cramer, H. (1938). Sur un nouveau theoreme-limite de la theorie des
  probabilites. \emph{Actualites Scientifiques et Industrielles}, 736,
  5-23.
\item
  Hoeffding, W. (1963). Probability inequalities for sums of bounded
  random variables. \emph{Journal of the American Statistical
  Association}, 58(301), 13-30.
\item
  Hoeffding, W. (1965). Asymptotically optimal tests for multinomial
  distributions. \emph{Annals of Mathematical Statistics}, 36(2),
  369-401.
\item
  Le Cam, L. (1986). \emph{Asymptotic Methods in Statistical Decision
  Theory}. Springer.
\item
  van der Vaart, A. W. (1998). \emph{Asymptotic Statistics}. Cambridge
  University Press.
\item
  Dembo, A., \& Zeitouni, O. (2010). \emph{Large Deviations Techniques
  and Applications} (2nd ed.). Springer.
\item
  Cover, T. M. \& Thomas, J. A. (2006). \emph{Elements of Information
  Theory} (2nd ed.). Wiley.
\item
  Pinsker, M. S. (1964). \emph{Information and Information Stability of
  Random Variables and Processes}. Holden-Day.
\item
  Fano, R. M. (1961). \emph{Transmission of Information: A Statistical
  Theory of Communications}. MIT Press.
\item
  Pollard, D. (1981). Strong consistency of k-means clustering.
  \emph{Annals of Statistics}, 9(1), 135-140.
\item
  Vershynin, R. (2018). \emph{High-Dimensional Probability}. Cambridge
  University Press.
\item
  Dawid, A. P., \& Skene, A. M. (1979). Maximum likelihood estimation of
  observer error-rates using the EM algorithm. \emph{JRSS Series C},
  28(1), 20-28.
\item
  Aloise, D., et al.~(2009). NP-hardness of Euclidean sum-of-squares
  clustering. \emph{Machine Learning}, 75(2), 245-248.
\item
  Ostrovsky, R., et al.~(2013). The effectiveness of Lloyd-type methods
  for the k-means problem. \emph{JACM}, 59(6), 1-22.
\item
  Kumar, A. \& Kannan, R. (2010). Clustering with spectral norm and the
  k-means algorithm. \emph{FOCS}, 299-308.
\item
  Arthur, D. \& Vassilvitskii, S. (2007). k-means++: The advantages of
  careful seeding. \emph{SODA}, 1027-1035.
\item
  von Luxburg, U. (2010). Clustering stability: An overview.
  \emph{Foundations and Trends in Machine Learning}, 2(3), 235-274.
\item
  Hubert, L., \& Arabie, P. (1985). Comparing partitions. \emph{Journal
  of Classification}, 2(1), 193-218.
\item
  Landis, J. R., \& Koch, G. G. (1977). The measurement of observer
  agreement for categorical data. \emph{Biometrics}, 33(1), 159-174.
\item
  Shamir, O., \& Tishby, N. (2009). On the reliability of clustering
  stability in the large sample regime. \emph{NIPS 2008}.
\item
  Rakhlin, A. \& Caponnetto, A. (2007). Stability of k-means clustering.
  \emph{NIPS 2007}.
\item
  Bousquet, O. (2002). A Bennett concentration inequality and its
  application to suprema of empirical processes. \emph{Comptes Rendus
  Mathematique}, 334(6), 495-500.
\item
  Shevtsova, I. G. (2011). On the absolute constants in the
  Berry-Esseen-type inequalities. \emph{Doklady Mathematics}, 83(3),
  320-323.
\item
  Jensen, J. L. (1995). \emph{Saddlepoint Approximations}. Oxford
  University Press.
\item
  Ingster, Y. I., \& Suslina, I. A. (2003). \emph{Nonparametric
  Goodness-of-Fit Testing Under Gaussian Models}. Springer.
\item
  Bartlett, P. L., \& Mendelson, S. (2002). Rademacher and Gaussian
  complexities. \emph{JMLR}, 3, 463-482.
\item
  Menon, A. K., et al.~(2015). Learning from corrupted binary labels via
  class-probability estimation. \emph{ICML}.
\item
  Northcutt, C. G., et al.~(2021). Confident learning: Estimating
  uncertainty in dataset labels. \emph{JAIR}, 70, 1373-1411.
\end{enumerate}

\begin{center}\rule{0.5\linewidth}{0.5pt}\end{center}

\emph{End of THEOREMS\_UNIFIED.md -- Single source of truth for the SCX
mathematical theory chain.}

\end{document}
